%% toc: P15 162 162.2,162.6,162.10
\section{Calibrating Wald's boundaries: proofs}
\label{sec:annex-p15-w02c-p15-020-calibration}

Throughout, histories of length \(m\) carry the product measure
\(\nu^m\), and under \(H_i\) the restriction of the law of the
observations to \(\mathcal F_m\) has density \(p_{im}\) with respect
to it.

\begin{annexitem}{annex:p15-w02c-p15-020-criteria-coincide}{thm:p15-w02c-p15-020-criteria-coincide}{Admissibility, the maximum and weighted sums}
\begin{proof}
Write \(e_i(S)=E_i(n\mid S)\) and \(e_i^*=e_i(S^*)\), which are finite.

(1) implies (2). Let \(S\) be admissible. By hypothesis
\(e_0^*\leq e_0(S)\) and \(e_1^*\leq e_1(S)\). If either inequality
were strict, \(S^*\) would be better than \(S\) and of equal strength,
contradicting admissibility. Hence both are equalities.

(2) implies (1). Let \(S\) satisfy (2) and suppose some
\(S'\in\mathcal S\) is better than \(S\). Then \(e_i(S')\leq e_i(S)=e_i^*\)
for both \(i\), with strict inequality for at least one \(i\). That
contradicts the minimality of \(e_i^*\) for that \(i\).

(2) implies (3). Such an \(S\) is admissible by the previous step, and
for every \(S'\in\mathcal S\),
\(\max\{e_0(S),e_1(S)\}=\max\{e_0^*,e_1^*\}\leq\max\{e_0(S'),e_1(S')\}\)
because \(e_i^*\leq e_i(S')\) for both \(i\).

(3) implies (1) by the definition of an optimum test.

(2) is equivalent to (4). For every \(S'\in\mathcal S\),
\[
 w_0e_0(S')+w_1e_1(S')-(w_0e_0^*+w_1e_1^*)
 =w_0(e_0(S')-e_0^*)+w_1(e_1(S')-e_1^*),
\]
a sum of two nonnegative terms, possibly infinite. So the minimum of the
weighted sum over \(\mathcal S\) is \(w_0e_0^*+w_1e_1^*\), attained by
\(S^*\). Since \(w_0,w_1>0\), a test attains it exactly when both terms
vanish, which is (2).

\(S^*\) satisfies (2) trivially, so it is an optimum test. The final
assertion of the statement is the arithmetic
\(\max\{e_0(S),e_1(S)\}=e_1(S)=e_1^*=\max\{e_0^*,e_1^*\}\), using
\(e_1^*\geq e_0(S)>e_0^*\).
\end{proof}
\end{annexitem}

\begin{annexitem}{annex:p15-w02c-p15-020-posterior-majority}{thm:p15-w02c-p15-020-posterior-majority}{Majority posteriors and a minority acceptance}
\begin{proof}
The posterior probabilities satisfy \(g_{0m}+g_{1m}=1\), as in the proof
of \cref{thm:p15-se-posterior}.

(1) If the rule accepts \(H_i\), then \(g_{im}\geq d_i>1/2\), so
\(g_{1-i,m}=1-g_{im}<1/2<g_{im}\).

(2) With the stated laws, \(f_0(1)=3/4\) and \(f_1(1)=1/2\). After a
first observation equal to \(1\),
\[
 g_{11}=\frac{\tfrac12\cdot\tfrac12}{\tfrac12\cdot\tfrac34+\tfrac12\cdot\tfrac12}
 =\frac{1/4}{5/8}=\frac25,\qquad g_{01}=\frac35 .
\]
Thus \(g_{11}\geq d_1=2/5\) and \(g_{01}<d_0=3/4\): only the condition
for \(H_1\) holds, so the rule accepts \(H_1\) at the first stage,
although \(g_{01}>g_{11}\). The first observation equals \(1\) with
probability \(f_0(1)=3/4\) under \(H_0\) and \(f_1(1)=1/2\) under
\(H_1\), so this happens with positive probability under either
hypothesis. Exchanging the roles of the two hypotheses and of the two
levels gives the case \(d_0\leq1/2\).
\end{proof}
\end{annexitem}

\begin{annexitem}{annex:p15-w02c-p15-020-fixed-size}{thm:p15-w02c-p15-020-fixed-size}{A stopped procedure is a test on the full sample}
\begin{proof}
Let \(U\) be the independent seed of \(\mathcal F_0\), with the same law
\(\lambda\) under both hypotheses; for a nonrandomized procedure take
\(U\) constant. On \(\{T=m\}\), \(m\leq N\), the decision \(D\) is a
measurable function of \(U\) and \(x_1,\ldots,x_m\), hence of \(U\)
and the full sample \(x=(x_1,\ldots,x_N)\). Since \(T\leq N\) almost
surely under both laws, there is a measurable \(d(u,x)\in\{0,1\}\)
with \(D=d(U,X)\) almost surely under both laws. Put
\[
 \psi(x)=\int\mathbf1_{\{d(u,x)=1\}}\,\lambda(du)\in[0,1].
\]
It is measurable by Tonelli's theorem, and by independence of \(U\)
and the sample and Fubini's theorem,
\(E_i\psi(X)=P_i(D=1)\). Hence \(E_0\psi=\alpha'\) and
\(E_1\psi=1-\beta'\), so \(\psi\) is a test on the \(N\)-observation
space with densities \(p_{0N},p_{1N}\) with respect to \(\nu^N\).

If \(0<\alpha'<1\), \cref{thm:p15-neyman-pearson} gives a most
powerful test \(\varphi^*\) of size \(\alpha'\) with
\(E_1\varphi^*\geq E_1\psi=1-\beta'\), so its second-kind error
\(1-E_1\varphi^*\) is at most \(\beta'\). If \(\alpha'=0\), the level
zero case of \cref{thm:p15-neyman-pearson} gives a level-zero test with
the maximum power, which is at least \(E_1\psi\). If \(\alpha'=1\),
the test identically equal to one has second-kind error \(0\leq\beta'\).
\end{proof}
\end{annexitem}

\begin{annexitem}{annex:p15-w02c-p15-020-one-boundary}{thm:p15-w02c-p15-020-one-boundary}{A pathwise comparison of nested continuation intervals}
\begin{proof}
Fix a sample path, with ratios \(L_m\in[0,\infty]\). Let \(T_j\) be the
first \(m\geq1\) with \(L_m\notin(B,A_j)\), and \(Q_1^{(j)},Q_0^{(j)}\)
the corresponding acceptance events.

Because \((B,A_1)\subseteq(B,A_2)\), any \(m\) with
\(L_m\notin(B,A_2)\) also has \(L_m\notin(B,A_1)\); so \(T_1\leq T_2\).

Suppose the path lies in \(Q_1^{(2)}\): \(T_2=m<\infty\),
\(L_m\geq A_2\). Then \(T_1=k\leq m\). If \(k<m\), then
\(L_k\in(B,A_2)\) because the second test had not stopped, and
\(L_k\notin(B,A_1)\); together these give \(L_k\geq A_1\). If \(k=m\),
then \(L_k\geq A_2\geq A_1\). In both cases the path is in
\(Q_1^{(1)}\).

Suppose the path lies in \(Q_0^{(1)}\): \(T_1=k<\infty\), \(L_k\leq B\),
and \(L_j\in(B,A_1)\subseteq(B,A_2)\) for \(j<k\). Then the second test
has not stopped before \(k\) and stops at \(k\) with \(L_k\leq B\), so
the path is in \(Q_0^{(2)}\).

Thus \(Q_1^{(2)}\subseteq Q_1^{(1)}\) and
\(Q_0^{(1)}\subseteq Q_0^{(2)}\), and for any law \(P\),
\(P(Q_1^{(2)})\leq P(Q_1^{(1)})\) and
\(P(Q_0^{(2)})\geq P(Q_0^{(1)})\). With \(P=P_0\) the first gives the
statement about \(\alpha'=P_0(Q_1)\); with \(P=P_1\) the second gives
the statement about \(\beta'=P_1(Q_0)\).

For the lower boundary let \(0<B_2\leq B_1<A\), now with \(T_j\) the
first \(m\) with \(L_m\notin(B_j,A)\) and acceptance events
\(Q_1^{(j)},Q_0^{(j)}\). Since \((B_1,A)\subseteq(B_2,A)\), the first
paragraph gives \(T_1\leq T_2\).
Suppose the path lies in \(Q_0^{(2)}\): \(T_2=m<\infty\),
\(L_m\leq B_2\). Then \(T_1=k\leq m\). If \(k<m\), then
\(L_k\in(B_2,A)\) and \(L_k\notin(B_1,A)\), so \(L_k\leq B_1\). If
\(k=m\), then \(L_k\leq B_2\leq B_1\). In both cases the path is in
\(Q_0^{(1)}\).
Suppose the path lies in \(Q_1^{(1)}\): \(T_1=k<\infty\), \(L_k\geq A\),
and \(L_j\in(B_1,A)\subseteq(B_2,A)\) for \(j<k\). Then the second test
stops at \(k\) with \(L_k\geq A\), so the path is in \(Q_1^{(2)}\).
Thus \(Q_0^{(2)}\subseteq Q_0^{(1)}\) and
\(Q_1^{(1)}\subseteq Q_1^{(2)}\). With \(P=P_1\) the first inclusion
shows that \(\beta'\) does not increase; with \(P=P_0\) the second shows
that \(\alpha'\) does not decrease.

Finally, if no value of \(L_m\) on the path lies in \([A_1,A_2)\), then
\(L_m\notin(B,A_1)\) and \(L_m\notin(B,A_2)\) are the same condition for
every \(m\), so the two tests stop and decide identically; the same
holds for \((B_2,B_1]\).
\end{proof}
\end{annexitem}

\begin{annexitem}{annex:p15-w02c-p15-020-nominal-sharp}{thm:p15-w02c-p15-020-nominal-sharp}{The equality case}
\begin{proof}
Here \(b(\alpha,\beta)<1<a(\alpha,\beta)\), since \(\alpha+\beta<1\).
By \cref{thm:p15-se-termination} the test terminates almost surely under
both laws, so \(P_1(Q_1)=1-\beta'\) and \(P_0(Q_0)=1-\alpha'\).

For \(m\geq1\), let \(H_m^1\) and \(H_m^0\) be the sets of length-\(m\)
histories on which the test stops at \(m\) and accepts \(H_1\),
respectively \(H_0\). On \(H_m^1\) the ratio is at least \(A>1\); the
convention that \(0/0=1\) excludes histories with \(p_{0m}=p_{1m}=0\),
so \(p_{1m}\geq Ap_{0m}\) there. On \(H_m^0\) the ratio is at most
\(B<1\), which excludes both \(p_{0m}=p_{1m}=0\) and \(p_{0m}=0<p_{1m}\),
so \(p_{1m}\leq Bp_{0m}\) there. Therefore
\[
 (1-\beta')-A\alpha'=\sum_{m\geq1}\int_{H_m^1}(p_{1m}-Ap_{0m})\,d\nu^m,
\]
\[
 B(1-\alpha')-\beta'=\sum_{m\geq1}\int_{H_m^0}(Bp_{0m}-p_{1m})\,d\nu^m,
\]
sums of integrals of nonnegative functions. This reproves the two
inequalities of \cref{thm:p15-se-errors} and identifies their equality
cases. Let \(D_m=H_m^1\cap\{p_{1m}\neq Ap_{0m}\}\), on which the first
integrand is strictly positive and \(p_{1m}>0\). The first sum is zero
iff \(\nu^m(D_m)=0\) for every \(m\), iff \(P_1(D_m)=\int_{D_m}p_{1m}\,d\nu^m=0\)
for every \(m\). Since on \(H_m^1\) the condition \(p_{1m}\neq Ap_{0m}\)
is exactly \(L_m\neq A\), this says \(P_1(Q_1\cap\{L_T\neq A\})=0\).
In the same way, with \(D_m'=H_m^0\cap\{p_{1m}\neq Bp_{0m}\}\), on which
\(p_{0m}>0\), the second sum is zero iff
\(P_0(Q_0\cap\{L_T\neq B\})=0\).

Now multiply \(A\alpha'\leq1-\beta'\) by \(\alpha>0\) and
\(\beta'\leq B(1-\alpha')\) by \(1-\alpha>0\). With the stated \(A,B\)
this gives
\[
 (1-\beta)\alpha'\leq\alpha(1-\beta'),\qquad
 (1-\alpha)\beta'\leq\beta(1-\alpha'),
\]
each with equality exactly when the corresponding unmultiplied
inequality is an equality. Adding and cancelling \(-\alpha\beta'\) and
\(-\beta\alpha'\) from both sides gives \(\alpha'+\beta'\leq\alpha+\beta\),
with equality iff both are equalities, that is, iff the test has no
overshoot at either boundary.

(1) If \(\alpha'>\alpha\) and \(\beta'>\beta\), then
\(\alpha'+\beta'>\alpha+\beta\), which is impossible.

(2) We have shown that equality of the sums is equivalent to no
overshoot. If \(\alpha'=\alpha\) and \(\beta'=\beta\), the sums are
equal. Conversely, if both multiplied inequalities are equalities,
\((\alpha',\beta')\) solves the linear system
\[
 (1-\beta)\alpha'+\alpha\beta'=\alpha,\qquad
 \beta\alpha'+(1-\alpha)\beta'=\beta,
\]
whose determinant \((1-\alpha)(1-\beta)-\alpha\beta=1-\alpha-\beta\) is
positive. Its unique solution is \((\alpha,\beta)\), which visibly
satisfies it.

(3) By \cref{thm:p15-se-nominal},
\(\alpha'\leq\alpha/(1-\beta)\), so
\(\alpha'-\alpha\leq\alpha/(1-\beta)-\alpha=\alpha\beta/(1-\beta)\);
the bound for \(\beta'\) is the same with the roles exchanged. At
\(\alpha=\beta=1/20\) both bounds equal
\((1/400)/(19/20)=1/380\).

The final statement follows from (2): if \((\alpha',\beta')\neq(\alpha,\beta)\),
the sums are not equal, and by the inequality already proved,
\(\alpha'+\beta'<\alpha+\beta\).
\end{proof}
\end{annexitem}

\begin{annexitem}{annex:p15-w02c-p15-020-exact-pair}{thm:p15-w02c-p15-020-exact-pair}{Widening}
\begin{proof}
The exact test has \(P_0\neq P_1\), so \cref{thm:p15-se-errors}
applies to it with actual errors \(\alpha,\beta\):
\(A^*\alpha\leq1-\beta\) and \(\beta\leq B^*(1-\alpha)\). Dividing,
\(A^*\leq(1-\beta)/\alpha=a(\alpha,\beta)\) and
\(B^*\geq\beta/(1-\alpha)=b(\alpha,\beta)\).
Apply \cref{thm:p15-w02c-p15-020-one-boundary} twice: first replace
\(A^*\) by \(a(\alpha,\beta)\) with lower boundary \(B^*\), then replace
\(B^*\) by \(b(\alpha,\beta)\) with upper boundary \(a(\alpha,\beta)\).
Each step requires lower below upper, which holds because
\(B^*<A^*\leq a(\alpha,\beta)\) and \(b(\alpha,\beta)\leq B^*\).
Neither step decreases the stopping time on any path, so the stopping
time \(T\) of the usual test dominates the stopping time \(T^*\) of the
exact one pathwise, and \(E_iT\geq E_iT^*\) for \(i=0,1\), in
\([0,\infty]\).

For the nonuniqueness remark, if two boundary pairs produce the same
stopping time and decision on every path, their acceptance events are
equal, so their errors are equal under every law. The last paragraph of
the proof of \cref{thm:p15-w02c-p15-020-one-boundary} gives such pairs
whenever no attainable ratio value lies between them.
\end{proof}
\end{annexitem}

\begin{annexitem}{annex:p15-w02c-p15-020-root-exists}{thm:p15-w02c-p15-020-root-exists}{Convexity, slope at zero and growth}
\begin{proof}
Since \(|z|\leq e^z+e^{-z}\), assumption (1) at \(t=\pm1\) gives
\(E|z|<\infty\).

\emph{Convexity and continuity.} For real \(s\neq t\) and
\(\lambda\in(0,1)\), convexity of the exponential gives
\(e^{(\lambda s+(1-\lambda)t)z}\leq\lambda e^{sz}+(1-\lambda)e^{tz}\),
with strict inequality wherever \(sz\neq tz\), that is, on
\(\{z\neq0\}\). By (2) this event has positive probability, so taking
expectations gives
\(\varphi(\lambda s+(1-\lambda)t)<\lambda\varphi(s)+(1-\lambda)\varphi(t)\):
\(\varphi\) is strictly convex on \(\mathbb R\). A finite convex function
on an open interval is continuous, so \(\varphi\) is continuous.
A strictly convex function takes any given value at most twice, since
three points with equal values would violate strict convexity at the
middle one.

\emph{Slope at zero.} For each fixed outcome, \(t\mapsto(e^{tz}-1)/t\) is
nondecreasing on \((0,\infty)\), being a difference quotient of the
convex function \(t\mapsto e^{tz}\), and it tends to \(z\) as
\(t\downarrow0\). For \(0<t\leq1\) it is bounded above by
\(e^z-1\), which is integrable. Monotone convergence applied to the
nonnegative variables \(e^z-1-(e^{tz}-1)/t\) increasing to
\(e^z-1-z\) gives
\[
 \frac{\varphi(t)-1}{t}\longrightarrow Ez\qquad(t\downarrow0).
\]

\emph{Growth.} Because \(\{z>0\}=\bigcup_{r\geq1}\{z\geq1/r\}\) is an
increasing union of positive probability, some \(\varepsilon>0\) has
\(P(z\geq\varepsilon)>0\), and then
\(\varphi(t)\geq e^{t\varepsilon}P(z\geq\varepsilon)\to\infty\) as
\(t\to\infty\). Likewise \(\varphi(t)\to\infty\) as \(t\to-\infty\).

\emph{Jensen.} The exponential is strictly convex and, for \(t\neq0\),
\(tz\) is not almost surely constant, because \(z\) takes positive and
negative values with positive probability. Hence
\(\varphi(t)>e^{tEz}\) for every \(t\neq0\).

Suppose \(Ez<0\). The slope at zero gives \(t_0>0\) with
\(\varphi(t_0)<1\), and growth gives \(t_1>t_0\) with \(\varphi(t_1)>1\).
By continuity some \(h\in(t_0,t_1)\) has \(\varphi(h)=1\). Since
\(\varphi(0)=1\) and \(\varphi\) takes the value one at most twice,
\(h\) is the only nonzero root. (Also, for \(t<0\),
\(\varphi(t)>e^{tEz}>1\), consistent with this.)

If \(Ez>0\), apply the previous case to \(-z\), which satisfies (1) and
(2) and has generating function \(t\mapsto\varphi(-t)\); its unique
nonzero root \(h'>0\) gives the unique nonzero root \(h=-h'<0\) of
\(\varphi\).

If \(Ez=0\), Jensen gives \(\varphi(t)>e^0=1\) for every \(t\neq0\).
\end{proof}
\end{annexitem}

\begin{annexitem}{annex:p15-w02c-p15-020-root-simple-laws}{thm:p15-w02c-p15-020-root-simple-laws}{Specified laws}
\begin{proof}
Equivalence gives \(P_1(f_0=0)=0\) because \(P_0(f_0=0)=0\), and
likewise \(P_0(f_1=0)=0\); so \(z\) is finite almost surely under both
laws. Under \(H_0\),
\[
 E_0e^{z}=\int_{\{f_0>0\}}\frac{f_1}{f_0}f_0\,d\nu
 =P_1(f_0>0)=1,
\]
and \(E_0e^{0\cdot z}=1\).

The set \(D=\{t:E_0e^{tz}<\infty\}\) is an interval: if \(s<t\) are in
\(D\) and \(\lambda\in(0,1)\), then
\(e^{(\lambda s+(1-\lambda)t)z}\leq\lambda e^{sz}+(1-\lambda)e^{tz}\).
It contains \(0\) and \(1\), hence \([0,1]\).

The score is not \(P_0\)-almost surely constant. If \(z=c\) almost
surely under \(P_0\), then \(f_1=e^cf_0\) \(\nu\)-almost everywhere on
\(\{f_0>0\}\); and \(f_1=0\) \(\nu\)-almost everywhere on \(\{f_0=0\}\),
because \(\int_{\{f_0=0\}}f_1\,d\nu=P_1(f_0=0)=0\). Integrating,
\(1=e^c\), so \(f_1=f_0\) \(\nu\)-almost everywhere and \(P_1=P_0\),
contrary to hypothesis. Hence \(z\neq0\) with positive \(P_0\)
probability, and the convexity argument in the proof of
\cref{thm:p15-w02c-p15-020-root-exists} shows that
\(t\mapsto E_0e^{tz}\) is strictly convex on \(D\). It equals one at
\(0\) and at \(1\), so it equals one at no third point of \(D\); outside
\(D\) it is infinite. Thus \(t=1\) is the only nonzero root.

Under \(H_1\), \(E_1e^{-z}=\int_{\{f_1>0\}}f_0\,d\nu=P_0(f_1>0)=1\).
The score is not \(P_1\)-almost surely constant, by the same argument
with the two laws exchanged, and the same convexity argument on the
interval of finiteness, which contains \([-1,0]\), shows that \(t=-1\)
is the only nonzero root.
\end{proof}
\end{annexitem}

\begin{annexitem}{annex:p15-w02c-p15-020-root-failure}{thm:p15-w02c-p15-020-root-failure}{A polynomial tail}
\begin{proof}
The function \(w\mapsto(1+w)^{-2}\) is continuous, decreasing from
\(1\) at \(w=0\) to \(0\), so it is the tail of a law on
\([0,\infty)\) without atoms. Then
\(EW=\int_0^\infty(1+w)^{-2}\,dw=1\), so \(Ez=-1\) and
\(E|z|\leq EW+2<\infty\). Next \(P(z>0)=P(W>2)=1/9\), and, as
\(P(W=2)=0\), \(P(z<0)=1-1/9=8/9\).

For \(t>0\), \(e^{tW}\geq(tW)^3/6\), and
\[
 EW^3=\int_0^\infty3w^2P(W>w)\,dw=\int_0^\infty\frac{3w^2}{(1+w)^2}\,dw=\infty,
\]
since the integrand tends to \(3\). Hence
\(\varphi(t)=e^{-2t}Ee^{tW}=\infty\).

For \(t<0\), \(tz\) is not almost surely constant, so strict Jensen
gives \(\varphi(t)>e^{tEz}=e^{-t}>1\). Hence \(\varphi(h)=1\) forces
\(h=0\).
\end{proof}
\end{annexitem}

\begin{annexitem}{annex:p15-w02c-p15-020-autoregressive}{thm:p15-w02c-p15-020-autoregressive}{Drift of the log ratio along a variance clock}
\begin{proof}
(1) In \(p_{i,m+1}\), the coordinate \(x_{m+1}\) appears only in the
factor
\[
 (2\pi)^{-1/2}\exp\{-\tfrac12(x_{m+1}-\lambda_ix_m)^2\},
\]
a
normal density in \(x_{m+1}\) with mean \(\lambda_ix_m\). Its integral
is one, leaving \(p_{im}\).

(2) The \(x_1\) factors cancel in the ratio. For real \(x,y\),
\[
 (x-\lambda_0y)^2-(x-\lambda_1y)^2
 =2xy(\lambda_1-\lambda_0)-y^2(\lambda_1^2-\lambda_0^2)
 =2\delta y(x-\bar\lambda y),
\]
which gives the displayed formula, with \(Z_m=\log(p_{1m}/p_{0m})\) and
\(Z_1=0\). Under \(H_i\) put \(e_j=x_j-\lambda_ix_{j-1}\) for
\(j\geq2\); these are independent standard normal variables, independent
of \(x_1\), and \(\mathcal F_{j-1}=\sigma(x_1,\ldots,x_{j-1})
=\sigma(x_1,e_2,\ldots,e_{j-1})\), so \(e_j\) is independent of
\(\mathcal F_{j-1}\). Since
\(x_j-\bar\lambda x_{j-1}=e_j+(\lambda_i-\bar\lambda)x_{j-1}\) and
\(\lambda_1-\bar\lambda=\delta/2=-(\lambda_0-\bar\lambda)\),
\[
 Z_m=M_m\pm\tfrac12\delta^2V_m,\qquad
 M_m=\delta\sum_{j=2}^mx_{j-1}e_j,\quad
 V_m=\sum_{j=2}^mx_{j-1}^2,
\]
with the plus sign under \(H_1\) and the minus sign under \(H_0\).

Each \(x_k\) is a linear combination of \(x_1,e_2,\ldots,e_k\), hence
normal with finite variance. The increment \(\delta x_{j-1}e_j\) is a
product of an \(\mathcal F_{j-1}\)-measurable square-integrable variable
and an independent standard normal one, so it is square integrable,
\(E(\delta x_{j-1}e_j\mid\mathcal F_{j-1})=0\), and
\(E((\delta x_{j-1}e_j)^2\mid\mathcal F_{j-1})=\delta^2x_{j-1}^2\).
With the shifted index \(X_n=M_{n+1}\) and filtration
\(\mathcal G_n=\mathcal F_{n+1}\), \(X\) is a square-integrable
martingale with \(X_0=0\) whose variance process, in the sense of
\cref{def:p15-predictable-variation}, is \(A_n=\delta^2V_{n+1}\).

\(V_\infty=\infty\) almost surely. On \(\{V_\infty<\infty\}\),
\(x_k\to0\), so \(e_k=x_k-\lambda_ix_{k-1}\to0\). But the independent
events \(\{|e_k|>1\}\), \(k\geq2\), have a common positive probability,
so by \cref{thm:p15-borel-cantelli-second} infinitely many occur almost
surely, and \(e_k\not\to0\) almost surely. Hence
\(P_i(V_\infty<\infty)=0\).

Apply \cref{thm:p15-variance-clock-normalization} with \(f(t)=1+t\),
for which \(\int_0^\infty f(t)^{-2}\,dt=1\): almost surely on
\(\{A_\infty=\infty\}\), which has probability one,
\(M_m/(1+\delta^2V_m)\to0\). Therefore
\[
 \frac{Z_m}{1+\delta^2V_m}\longrightarrow\pm\frac12
\]
almost surely, and as \(V_m\to\infty\), \(Z_m\to+\infty\) under \(H_1\)
and \(Z_m\to-\infty\) under \(H_0\).

(3) Since \(Z_m\) tends to \(\pm\infty\), it eventually leaves the
bounded interval \((\log B,\log A)\) almost surely under each
hypothesis; so the test stops at a finite time almost surely under both.
The densities \(p_{im}\) are the joint densities of the actual
histories, consistent by (1). \Cref{thm:p15-se-dependent} therefore
gives \(A\alpha'\leq1-\beta'\) and \(\beta'\leq B(1-\alpha')\).
When \(\alpha',\beta'<1\), dividing by \(1-\beta'\) and \(1-\alpha'\)
gives the source's quotient form.
\end{proof}
\end{annexitem}

\begin{annexitem}{annex:p15-w02c-p15-020-gamma-bar}{thm:p15-w02c-p15-020-gamma-bar}{Two-sided overshoot of at most one score}
\begin{proof}
Let \(h>0\). By \cref{thm:p15-se-oc}, \(Ee^{hZ_T}=1\). Before \(T\) the
sum lies in \((b,a)\) and the last step has \(|z_T|\leq c\), so
\(a\leq Z_T\leq a+c\) on \(Q_1\) and \(b-c\leq Z_T\leq b\) on \(Q_0\).
Hence \(A^h\leq e^{hZ_T}\leq sA^h\) on \(Q_1\) and
\(B^h/s\leq e^{hZ_T}\leq B^h\) on \(Q_0\). Since \(T<\infty\) almost
surely, \(Q_0\) and \(Q_1\) partition the sample space up to a null
set, and taking expectations of the two bounds gives
\[
 \gamma A^h+(1-\gamma)B^h/s\leq1\leq\gamma sA^h+(1-\gamma)B^h .
\]
As \(A^h>1>B^h\) and \(s\geq1\), the coefficients \(A^h-B^h/s\) and
\(sA^h-B^h\) are positive, and solving for \(\gamma\) gives the two
stated bounds.

Let \(h<0\). Then \(z'=-z\) has \(Ee^{(-h)z'}=1\) with \(-h>0\), the
same bound \(c\), and the test for \(z'\) has upper and lower log
boundaries \(-b\) and \(-a\), that is \(A'=1/B\) and \(B'=1/A\). Its
upper exits are the lower exits of the original walk, so
\(\gamma'=1-\gamma\). Since \(A'^{-h}=B^h\) and \(B'^{-h}=A^h\), the
case already proved gives the bounds for \(1-\gamma\) with \(A^h\) and
\(B^h\) exchanged, and their common value at \(s=1\) is
\((1-A^h)/(B^h-A^h)=1-\bar\gamma\).

For uniformity, let \(K\subset(0,\infty)\) be compact. The two bounds
are continuous functions of \((h,s)\) on \(K\times[1,2]\), because
their denominators are positive there, and they equal \(\bar\gamma\)
at \(s=1\). A continuous function on a compact set is uniformly
continuous, so the bounds converge to \(\bar\gamma\) uniformly in
\(h\in K\) as \(s\downarrow1\). A compact subset of
\(\mathbb R\setminus\{0\}\) is the union of a compact subset of
\((0,\infty)\) and one of \((-\infty,0)\), and the negative part is
handled by the reflection above. If \(|h_k|\leq H\) and \(c_k\to0\),
then \(s_k=e^{|h_k|c_k}\leq e^{Hc_k}\to1\), and since \(\gamma_k\) and
\(\bar\gamma_k\) lie between the same two bounds,
\(|\gamma_k-\bar\gamma_k|\) is at most their difference, which tends
to zero. Finally \(|Ez_k|\leq c_k\) and \(\operatorname{Var}z_k\leq Ez_k^2\leq c_k^2\).
For the specified hypotheses the roots are \(\pm1\)
(\cref{thm:p15-w02c-p15-020-root-simple-laws}), a compact set.
\end{proof}
\end{annexitem}

\begin{annexitem}{annex:p15-w02c-p15-020-gamma-bar-failure}{thm:p15-w02c-p15-020-gamma-bar-failure}{Rare jumps}
\begin{proof}
For \(0<p<1/e\), \(1-pe>0\) and \(1-p>1-pe\), so \(\epsilon_p>0\). Then
\(Ee^z=pe+(1-p)(1-pe)/(1-p)=1\). The score is bounded and takes both
signs, so by \cref{thm:p15-w02c-p15-020-root-exists} it has exactly one
nonzero root, which is therefore \(h=1\). By strict convexity of the
exponential, \(e^{Ez}<Ee^z=1\), so \(Ez<0\). As \(p\to0\),
\(\epsilon_p\to0\), so \(|Ez|\leq p+\epsilon_p\to0\) and
\(\operatorname{Var}z\leq Ez^2\leq p+\epsilon_p^2\to0\).

From any value \(x\in(b,a)=(-1/4,1/4)\), a step \(+1\) leads to
\(x+1>3/4>a\), so each step \(+1\) taken before exit ends the test with
an upper exit. Hence \(P(T>m)\leq(1-p)^m\) and \(T<\infty\) almost
surely. Steps \(-\epsilon_p\) decrease the sum, so an upper exit occurs
only through a step \(+1\) from some \(x>-1/4\), and then
\(Z_T-a=x+1-1/4>1/2\). A lower exit occurs through a step
\(-\epsilon_p\) from some \(x>b\), so \(b-\epsilon_p<Z_T\leq b\). By
\cref{thm:p15-se-oc}, \(Ee^{Z_T}=1\); bounding \(e^{Z_T}\) below by
\(Ae^{1/2}\) on \(Q_1\) and by \(Be^{-\epsilon_p}\) on \(Q_0\),
\[
 1\geq\gamma Ae^{1/2}+(1-\gamma)Be^{-\epsilon_p},
\]
which gives the stated bound on \(\gamma\), as \(Ae^{1/2}>1>Be^{-\epsilon_p}\).
Its limit follows from \(\epsilon_p\to0\), and the strict inequality
from \(Ae^{1/2}>A\) and \(1-B>0\). The value
\(\bar\gamma=(1-B)/(A-B)\) for \(h=1\) does not depend on \(p\), so
\(\liminf_{p\to0}(\bar\gamma-\gamma)\geq(1-B)/(A-B)-(1-B)/(Ae^{1/2}-B)>0\).
\end{proof}
\end{annexitem}

\begin{annexitem}{annex:p15-w02c-p15-020-lattice}{thm:p15-w02c-p15-020-lattice}{Monotone coupling}
\begin{proof}
(1) Since \(z^-\leq z\leq z^+\) observation by observation,
\(Z^-_m\leq Z_m\leq Z^+_m\) for every \(m\). Let \(Y\leq Y'\) be two
paths with \(Y\) leaving \((b,a)\) upward at time \(t\): \(b<Y_m<a\) for
\(m<t\) and \(Y_t\geq a\). Then \(Y'_m\geq Y_m>b\) for \(m<t\) and
\(Y'_t\geq a\), so \(Y'\) never reaches \((-\infty,b]\) up to time
\(t\) and reaches \([a,\infty)\) at some time at most \(t\): it also
leaves upward. Applied to the pairs \((Z^-,Z)\) and \((Z,Z^+)\), the
event of an upper exit of the smaller path is contained in that of the
larger, which gives the two inequalities. If \(|z|\leq c\), then
\(z^\pm/d\) is an integer of absolute value at most \(c/d+1\), so it
takes finitely many values, and the lattice walk with boundaries
\(a/d\) and \(b/d\) is of the kind treated in
\cref{thm:p15-se-finite-recursion}.

(2) If the law of \(z\) has no atoms, then for \(m\geq1\),
\(Z_m=z_1+Z'\) with \(Z'=z_2+\cdots+z_m\) independent of \(z_1\), and by
Fubini \(P(Z_m=a)=E\,P(z_1=a-y)|_{y=Z'}=0\); similarly for \(b\).
Fix a path with \(T=t<\infty\) and \(Z_m\notin\{a,b\}\) for all \(m\);
these paths have probability one. Then \(b<Z_m<a\) for \(m<t\), and
\(Z_t>a\) or \(Z_t<b\), all strictly. Let \(\epsilon>0\) be the smallest
distance from \(Z_1,\ldots,Z_t\) to \(\{a,b\}\). Since
\(|Z^\pm_m-Z_m|\leq md\leq td\) for \(m\leq t\), every \(d<\epsilon/t\)
gives the same strict inequalities for \(Z^\pm_1,\ldots,Z^\pm_t\), so
\(T^\pm=t\) with the same decision. Thus the indicator of an upper exit
of \(Z^\pm\) converges almost surely to that of \(Z\), and bounded
convergence gives \(\gamma^\pm_d\to\gamma\).

(3) On \(\{T^\pm<\infty\}\), \(Z^\pm_{T^\pm-1}\in(b,a)\), or
\(T^\pm=1\) and the previous value is zero, and the last step has
absolute value at most \(c+d\); so for \(d\leq1\),
\(|Z^\pm_{T^\pm}|\mathbf1_{\{T^\pm<\infty\}}\leq\max\{a,-b\}+c+1\). By
(2), almost surely \(T^\pm=T\) for small \(d\), and then
\(Z^\pm_{T^\pm}=Z^\pm_T\to Z_T\) because \(|Z^\pm_T-Z_T|\leq Td\).
Dominated convergence gives the convergence of the terminal means. When
\(Ez\neq0\) and \(ET<\infty\), \(EZ_T=Ez\cdot ET\) by
\cref{thm:p15-se-mean}, which gives the last assertion.
\end{proof}
\end{annexitem}

\begin{annexitem}{annex:p15-w02c-p15-020-small-shift}{thm:p15-w02c-p15-020-small-shift}{Limits of the exact identities as the shift vanishes}
\begin{proof}
Translate the observations by \((\theta_0+\theta_1)/2\). The score is
unchanged, \(z=(\theta_1-\theta_0)x-(\theta_1^2-\theta_0^2)/2=2d\tilde x\),
and \(\tilde x\) is \(N(-d,1)\) under \(H_0\) and \(N(d,1)\) under
\(H_1\). This is the setting of \cref{thm:p15-se-normal} with third law
\(N(\theta,1)\), \(\theta=\mp d\). There \(Ez=\mp2d^2\), and the test
terminates with finite expectation under both hypotheses by
\cref{thm:p15-se-termination}.

(1) By \cref{thm:p15-se-normal}, at \(\theta=d\),
\(\xi=2d\{d+\phi(d)/\Phi(d)\}\) and \(-\xi'=2d\{\phi(d)/(1-\Phi(d))-d\}\);
at \(\theta=-d\), using \(\phi(-d)=\phi(d)\) and
\(\Phi(-d)=1-\Phi(d)\), \(\xi=2d\{\phi(d)/(1-\Phi(d))-d\}\) and
\(-\xi'=2d\{\phi(d)/\Phi(d)+d\}\). Since \(\Phi(d)\geq1/2\geq1-\Phi(d)\),
each is at most \(M_d\). \Cref{thm:p15-se-mean} gives, with
\(\gamma=P(Q_1)\) under the law in force,
\[
 (1-\gamma)\xi'\leq Ez\,ET-\{(1-\gamma)b+\gamma a\}\leq\gamma\xi,
\]
so the middle term has absolute value at most \(M_d\). Under \(H_1\),
\(Ez=2d^2\) and \(\gamma=1-\beta'_d\); under \(H_0\), \(Ez=-2d^2\) and
\(\gamma=\alpha'_d\). This is the displayed pair. For \(0<d\leq1\),
\(\phi(d)\leq\phi(0)\) and \(1-\Phi(d)\geq1-\Phi(1)\), so
\(M_d\leq2d\{1+\phi(0)/(1-\Phi(1))\}\), and
\(\phi(0)/(1-\Phi(1))<0.40/0.158<2.6\).

(2) Under \(H_0\), \(\varphi(t)=\exp\{2d^2(t^2-t)\}\), so \(h=1\); under
\(H_1\), \(\varphi(t)=\exp\{2d^2(t^2+t)\}\), so \(h=-1\). In both cases
\cref{thm:p15-se-normal} gives \(\delta=\Phi(d)/(1-\Phi(d))\geq1\) and
\(\eta=1/\delta\), so \cref{thm:p15-se-oc} applies. Under \(H_0\) it
gives
\[
 \frac{1-B}{\delta A-B}\leq\alpha'_d\leq\frac{1-\eta B}{A-\eta B},
\]
and as \(d\to0\), \(\delta,\eta\to1\), so both bounds tend to
\[
 \frac{1-B}{A-B}
 =\frac{(1-\alpha-\beta)/(1-\alpha)}{\{(1-\beta)(1-\alpha)-\alpha\beta\}/\{\alpha(1-\alpha)\}}
 =\alpha .
\]
Under \(H_1\) the case \(h<0\) gives
\((1-A^{-1})/(\delta B^{-1}-A^{-1})\leq\beta'_d\leq(1-\eta A^{-1})/(B^{-1}-\eta A^{-1})\),
and both bounds tend to
\((1-A^{-1})/(B^{-1}-A^{-1})
=\{(1-\alpha-\beta)/(1-\beta)\}\big/\{(1-\alpha-\beta)/(\beta(1-\beta))\}=\beta\).

(3) As \(d\to0\), \(M_d\to0\), so by (1) and (2)
\(2d^2E_{1,d}T\to(1-\beta)a+\beta b\) and
\(2d^2E_{0,d}T\to-\{(1-\alpha)b+\alpha a\}\). Since
\(2d^2n_d=(q_{1-\alpha}+q_{1-\beta})^2/2\), dividing gives the limits
\(R_1\) and \(R_0\). Finally \(n_d\to\infty\) and
\(n_d\leq N_d<n_d+1\), so \(N_d/n_d\to1\).

(4) The laws are equivalent with integrable scores, \(\mu_0<0<\mu_1\),
and the test terminates, so \cref{thm:p15-se-information} gives
\(|\mu_i|E_{i,d}T\geq I_{i,d}\). With \(u=\alpha'_d\), \(v=\beta'_d\),
\[
 I_{1,d}=v\log\frac v{1-u}+(1-v)\log\frac{1-v}u,\qquad
 I_{0,d}=(1-u)\log\frac{1-u}v+u\log\frac u{1-v},
\]
continuous in \((u,v)\in(0,1)^2\). By (2) they tend to
\(\beta b+(1-\beta)a\) and \(-\{(1-\alpha)b+\alpha a\}\), which by (3)
are also the limits of \(|\mu_1|E_{1,d}T=2d^2E_{1,d}T\) and
\(|\mu_0|E_{0,d}T\). These limits are positive: \(I_{1,d}\) at
\((\alpha,\beta)\) is the relative entropy of the two-point law
\((\beta,1-\beta)\) with respect to \((1-\alpha,\alpha)\), which is
positive since \(\beta\neq1-\alpha\), and similarly for \(I_{0,d}\).
Hence the ratios tend to one.

For the example after the theorem, with \(p_0=(e^{2d}-1)/(e^{3d}-1)\in(0,1)\)
and \(p_1=e^dp_0=(e^{3d}-e^d)/(e^{3d}-1)\in(0,1)\),
\(f_1(1)/f_0(1)=e^d\) and
\[
 \frac{1-p_1}{1-p_0}=\frac{e^d-1}{e^{3d}-e^{2d}}=e^{-2d},
\]
so \(z\in\{d,-2d\}\). The two laws are equivalent, the scores bounded,
and \(E_0z<0<E_1z\) by the strict Jensen argument used above. With
\(A=e^d\) and \(B=e^{-2d}\), the first observation gives \(Z_1=\log A\)
or \(Z_1=\log B\), so \(T=1\) with no overshoot, and
\cref{thm:p15-se-no-overshoot} applies.
\end{proof}
\end{annexitem}

\begin{annexitem}{annex:p15-w02c-p15-020-moment-system}{thm:p15-w02c-p15-020-moment-system}{Order-by-order solution of the differentiated identity}
\begin{proof}
On the open interval where \(\varphi\) is finite it is infinitely
differentiable and strictly convex, since \(z\) is nonconstant. As
\(\varphi(0)=\varphi(h)=1\), its minimum lies strictly between \(0\)
and \(h\) and \(\varphi'\) is strictly increasing, so \(\varphi'(h)\)
is nonzero, with the sign of \(h\). Also \(\ell(0)=\ell(h)=0\),
\(\ell'(0)=-\mu\) and \(\ell'(h)=-\mu_h\). Put
\(F(t)=\bar\gamma A^t\widehat\Psi_+(\ell(t))+(1-\bar\gamma)B^t\widehat\Psi_-(\ell(t))\).

Order zero. \(F(0)=F(h)=1\) reads
\(\bar\gamma\widehat\Psi_+(0)+(1-\bar\gamma)\widehat\Psi_-(0)=1\) and
\(\bar\gamma A^h\widehat\Psi_+(0)+(1-\bar\gamma)B^h\widehat\Psi_-(0)=1\).
Its determinant is \(\bar\gamma(1-\bar\gamma)(B^h-A^h)\neq0\), and
\(\widehat\Psi_\pm(0)=1\) solves it, because
\(\bar\gamma A^h+(1-\bar\gamma)B^h=B^h+\bar\gamma(A^h-B^h)=1\).

Order \(j\geq1\). By the Leibniz rule,
\[
 \frac{d^j}{dt^j}\{A^tg(t)\}=\sum_{i=0}^j\binom ji(\log A)^{j-i}A^tg^{(i)}(t),
\]
and by the chain rule applied repeatedly, \(g=\widehat\Psi_+\circ\ell\)
has
\[
 g^{(i)}(t)=\widehat\Psi_+^{(i)}(\ell(t))\ell'(t)^i+P_i(t),
\]
where
\(P_i\) is a sum of products of derivatives of \(\ell\) with
\(\widehat\Psi_+^{(k)}(\ell(t))\), \(k<i\). So in \(F^{(j)}(0)=0\) the
only terms of order \(j\) are
\(\bar\gamma(-\mu)^j\widehat\Psi_+^{(j)}(0)+(1-\bar\gamma)(-\mu)^j\widehat\Psi_-^{(j)}(0)\),
and in \(F^{(j)}(h)=0\) they are
\(\bar\gamma A^h(-\mu_h)^j\widehat\Psi_+^{(j)}(0)+(1-\bar\gamma)B^h(-\mu_h)^j\widehat\Psi_-^{(j)}(0)\);
everything else involves orders below \(j\), which are known by
induction. The determinant of this \(2\times2\) system is
\(\bar\gamma(1-\bar\gamma)(-\mu)^j(-\mu_h)^j(B^h-A^h)
=\bar\gamma(1-\bar\gamma)\mu^j\mu_h^j(B^h-A^h)\), nonzero since
\(0<\bar\gamma<1\), \(\mu,\mu_h\neq0\), and \(A^h\neq B^h\) for
\(A>B\), \(h\neq0\). So each order has exactly one solution, and the
\(2r\) equations of orders \(1,\ldots,r\) determine the \(2r\)
unknowns.

For \(j=1\), with \(\widehat\Psi_\pm(0)=1\), \(F'(0)=0\) is
\(\bar\gamma(a-\mu\widehat\Psi_+'(0))+(1-\bar\gamma)(b-\mu\widehat\Psi_-'(0))=0\),
the first displayed equation, and \(F'(h)=0\) is
\(\bar\gamma A^h(a-\mu_h\widehat\Psi_+'(0))+(1-\bar\gamma)B^h(b-\mu_h\widehat\Psi_-'(0))=0\),
the second. Substituting \(v=S/\mu-u\), \(S=\bar\gamma a+(1-\bar\gamma)b\),
into the second gives \(\mu_h(A^h-B^h)u=S_h-\mu_hB^hS/\mu\) with
\(S_h=\bar\gamma aA^h+(1-\bar\gamma)bB^h\), which is the displayed
\(u\).
\end{proof}
\end{annexitem}

\begin{annexitem}{annex:p15-w02c-p15-020-one-sided}{thm:p15-w02c-p15-020-one-sided}{Reflection of the score and agreement on upper exits}
\begin{proof}
(1) \(Z_m\leq b\) if and only if \(-Z_m\geq-b=\log(1/B)\), so the two
stopping rules stop at the same time on every path. The score \(-z\)
has mean \(-Ez>0\) and variance \(\sigma_z^2\), and
\((-Ez)^2=(Ez)^2\), so \(m\) is the same function of \(n\). The upper
one-sided approximation for \(-z\) with upper log boundary \(-b\) has
parameter \((-Ez)(-b)/\sigma_z^2=(Ez/\sigma_z^2)\log B\), which is
positive since \(Ez<0\) and \(\log B<0\) for \(0<B<1\). The properties
of \(f_c\) for \(c>0\) are \cref{thm:p15-se-continuous-density}.

(2) On \(Q_1\), \(b<Z_m<a\) for \(m<T\) and \(Z_T\geq a\), so \(T\) is
the first time the sum reaches \([a,\infty)\), that is \(T=T_A\). Hence
\(\{T\neq T_A\}\subset Q_1^c\), whose probability is \(1-\gamma\). For
any set \(G\),
\(|P(T\in G)-P(T_A\in G)|=|P(T\in G,T\neq T_A)-P(T_A\in G,T\neq T_A)|\leq P(T\neq T_A)\).
\end{proof}
\end{annexitem}

\begin{annexitem}{annex:p15-w02c-p15-020-current-size}{thm:p15-w02c-p15-020-current-size}{Monotonicity and an affinity bound}
\begin{proof}
(1) A test based on \(x_1,\ldots,x_N\) is also a test based on
\(x_1,\ldots,x_{N+1}\) that ignores the last observation, with the same
errors. The infimum over the larger class is no larger.

(2) Let \(k=1/\alpha\) and consider the test that rejects \(H_0\) when
\(p_{1N}\geq kp_{0N}\). Its size is
\(P_0(p_{1N}\geq kp_{0N})\leq k^{-1}E_0(p_{1N}/p_{0N})\leq\alpha\), since
the set \(\{p_{0N}=0\}\) is \(P_0\)-null and
\(E_0(p_{1N}/p_{0N})=\int_{\{p_{0N}>0\}}p_{1N}\,d\nu^N\leq1\). On the
event \(\{p_{1N}<kp_{0N}\}\cap\{p_{1N}>0\}\), the indicator is at most
\(\sqrt{kp_{0N}/p_{1N}}\), and \(\{p_{1N}=0\}\) is \(P_1\)-null, so the
second-kind error is at most
\[
 E_1\sqrt{kp_{0N}/p_{1N}}=\sqrt k\int\sqrt{p_{0N}p_{1N}}\,d\nu^N=\sqrt k\,\rho^N,
\]
the last step by Fubini, since the integrand is a product over the
coordinates. Hence \(\beta_N(\alpha)\leq\alpha^{-1/2}\rho^N\). By the
Cauchy--Schwarz inequality \(\rho\leq(\int f_0)^{1/2}(\int f_1)^{1/2}=1\),
with equality only if \(\sqrt{f_0}\) and \(\sqrt{f_1}\) are proportional
almost everywhere, which forces \(f_0=f_1\) almost everywhere and
\(P_0=P_1\). So \(\rho<1\).

(3) By (2) some \(N\) has \(\beta_N(\alpha)\leq\beta\); by (1) every
larger \(N\) does too. The least such \(N\) exists in the positive
integers.

For the normal hypotheses, after the translation used for
\cref{thm:p15-w02c-p15-020-small-shift}, the ratio test rejects when
\(\tilde x_1+\cdots+\tilde x_N\geq\kappa\), a sum that is
\(N(\mp Nd,N)\); size \(\alpha\) gives \(\kappa=-Nd+\sqrt Nq_{1-\alpha}\),
and the second-kind error is
\(\Phi((\kappa-Nd)/\sqrt N)=\Phi(q_{1-\alpha}-2d\sqrt N)\). This is at
most \(\beta=\Phi(-q_{1-\beta})\) exactly when
\(2d\sqrt N\geq q_{1-\alpha}+q_{1-\beta}\), that is \(N\geq n_d\).
\end{proof}
\end{annexitem}

\begin{annexitem}{annex:p15-w02c-p15-020-current-efficiency}{thm:p15-w02c-p15-020-current-efficiency}{A randomized competitor}
\begin{proof}
Randomization. Let a terminating procedure have errors \((u,v)\) with
\(u+v<1\), and let \(u\leq u^*\), \(v\leq v^*\), \(u^*+v^*\leq1\). Add
an independent seed and, at termination, change a decision \(0\) into
\(1\) with probability \(p\) and a decision \(1\) into \(0\) with
probability \(r\). This is a procedure in the sense of
\cref{def:p15-se-procedure} with the same stopping time, and its errors
are \(u+p(1-u)-ru\) and \(v+r(1-v)-pv\). Setting them equal to
\(u^*,v^*\) gives a linear system with determinant \(1-u-v>0\) and
solution
\[
 p=\frac{(u^*-u)(1-v)+(v^*-v)u}{1-u-v},\qquad
 r=\frac{(v^*-v)(1-u)+(u^*-u)v}{1-u-v},
\]
both nonnegative. Moreover \(p\leq1\) is equivalent to
\(u^*(1-v)+v^*u\leq1-v\), and indeed
\(u^*(1-v)+v^*u\leq(1-v^*)(1-v)+v^*u=1-v-v^*(1-u-v)\leq1-v\); by
symmetry \(r\leq1\).

The current test with \(N_d\) observations has errors
\((\alpha,\beta_{N_d}(\alpha))\), and its worst expected number of
observations is \(N_d\). By \cref{thm:p15-w02c-p15-020-current-size},
\(\beta_{N_d}(\alpha)=\Phi(q_{1-\alpha}-2d\sqrt{N_d})\), and
\(q_{1-\alpha}+q_{1-\beta}\leq2d\sqrt{N_d}<\sqrt{(q_{1-\alpha}+q_{1-\beta})^2+4d^2}\),
so \(\beta_{N_d}(\alpha)\leq\beta\) and \(\beta_{N_d}(\alpha)\to\beta\).

Let \(0<\tilde\alpha<\alpha\) and \(0<\tilde\beta<\beta\), and let
\(\tilde S_d\) be the ratio test with boundaries
\(a(\tilde\alpha,\tilde\beta)\), \(b(\tilde\alpha,\tilde\beta)\). By
\cref{thm:p15-w02c-p15-020-small-shift}, (2), its errors tend to
\((\tilde\alpha,\tilde\beta)\), so for all small \(d\) they are below
\(\alpha\) and \(\beta_{N_d}(\alpha)\) respectively and sum to less than
one. The randomization above, with
\((u^*,v^*)=(\alpha,\beta_{N_d}(\alpha))\), \(u^*+v^*\leq\alpha+\beta<1\),
turns \(\tilde S_d\) into a test of the same strength as the current
test with the same expected numbers of observations. By the definition
of minimax efficiency,
\[
 e_d\leq\frac{\max_iE_i\tilde T_d}{N_d}.
\]
By \cref{thm:p15-w02c-p15-020-small-shift}, (1) and (2) applied to
\((\tilde\alpha,\tilde\beta)\), \(2d^2E_1\tilde T_d\to G_1(\tilde\alpha,\tilde\beta)
=(1-\tilde\beta)\log\frac{1-\tilde\beta}{\tilde\alpha}+\tilde\beta\log\frac{\tilde\beta}{1-\tilde\alpha}\)
and \(2d^2E_0\tilde T_d\to G_0(\tilde\alpha,\tilde\beta)
=(1-\tilde\alpha)\log\frac{1-\tilde\alpha}{\tilde\beta}+\tilde\alpha\log\frac{\tilde\alpha}{1-\tilde\beta}\),
while \(2d^2N_d\to(q_{1-\alpha}+q_{1-\beta})^2/2\). Hence
\[
 \limsup_{d\to0}e_d\leq\frac{2\max_iG_i(\tilde\alpha,\tilde\beta)}{(q_{1-\alpha}+q_{1-\beta})^2}.
\]
The functions \(G_i\) are continuous on \((0,1)^2\), and
\(2G_i(\alpha,\beta)/(q_{1-\alpha}+q_{1-\beta})^2=R_i(\alpha,\beta)\);
letting \(\tilde\alpha\uparrow\alpha\) and \(\tilde\beta\uparrow\beta\)
proves the bound. The numerical values are evaluations of
\(R_0,R_1\) at the stated pairs.
\end{proof}
\end{annexitem}

\begin{annexitem}{annex:p15-w02c-p15-020-truncation-exact}{thm:p15-w02c-p15-020-truncation-exact}{Comparing decisions path by path}
\begin{proof}
The truncated and untruncated procedures make the same decision on
\(\{T\leq n_0\}\). On \(\{T>n_0\}\) the truncated one accepts \(H_1\)
exactly when \(Z_{n_0}>0\). Since \(T<\infty\) almost surely, \(Q_0\)
and \(Q_1\) partition the sample space up to a null set. Under \(H_0\),
\[
 \alpha(n_0)-\alpha'
 =P_0(T>n_0,Z_{n_0}>0)-P_0(T>n_0,Q_1),
\]
and splitting the first event by \(Q_0,Q_1\) and the second by the sign
of \(Z_{n_0}\), the common part \(P_0(T>n_0,Z_{n_0}>0,Q_1)\) cancels,
leaving \(\rho_0(n_0)-\rho_0^-(n_0)\). Under \(H_1\), the same argument
with the roles of the decisions exchanged gives the formula for
\(\beta(n_0)\).

On \(\{T>n_0\}\), \(b<Z_{n_0}<a\). Hence
\(\bar\rho_0(n_0)=P_0(T\leq n_0,0<Z_{n_0}<a)+P_0(T>n_0,Z_{n_0}>0,Q_0)+P_0(T>n_0,Z_{n_0}>0,Q_1)\),
so
\(\alpha+\bar\rho_0(n_0)-\alpha(n_0)=(\alpha-\alpha')+\rho_0^-(n_0)
+P_0(T\leq n_0,0<Z_{n_0}<a)+P_0(T>n_0,Z_{n_0}>0,Q_1)\). The
\(\beta\) case is the same with \(P_1\), the event \(b<Z_{n_0}\leq0\)
and the decisions exchanged.
\end{proof}
\end{annexitem}
